\Legend
小常和小鹏在一张有 $2^n$ 个点、$m$ 条有向边的有向无环图上玩游戏，点的编号为 $0$ 到 $2^n-1$。棋盘上有 $k$ 枚棋子，第 $i$ 枚初始位于点 $a_i$，多枚棋子可以位于同一个点。小常先手，双方轮流选择任意一枚棋子，沿当前点的一条出边移动它。轮到某人时，如果所有棋子都无法移动，则此人输。

两人首先共同决定初始位置 $a_i$。把每个 $a_i$ 写成恰好 $n$ 位的二进制数，最低位排在前面。按棋子顺序将这些位拼接成 $k n$ 个位置：先决定 $a_1$ 的第 $0$ 位，再决定第 $1$ 位，直至最高位；接着决定 $a_2$，依此类推。小常决定第 $1,3,5,\ldots$ 个位置上的 $0$ 或 $1$，小鹏决定第 $2,4,6,\ldots$ 个位置。双方都知道此前已决定的全部位，并以最终赢得棋子游戏为目标，采取最优策略。

求从决定初始位置开始，小常能否保证最终获胜。能则输出 \tt{First}，否则输出 \tt{Second}。

\Input

第一行三个整数 $n,m,k$。接下来 $m$ 行，每行两个整数 $u,v$，表示一条从 $u$ 到 $v$ 的有向边。

$1\le n\le14$，$1\le m\le 8\times10^4$，$1\le k\le20$；$0\le u,v<2^n$。输入保证图是有向无环图。

\Output

输出 \tt{First} 或 \tt{Second}。

\Examples

\begin{example}
\exmp{2 2 2
1 0
3 2}{First}
\end{example}

\begin{example}
\exmp{2 2 2
3 1
2 0}{Second}
\end{example}

\begin{example}
\exmp{3 6 3
7 3
6 2
5 1
4 0
7 5
6 4}{First}
\end{example}
